\documentclass[a4paper,11pt]{article}
\usepackage[T1]{fontenc}
\usepackage[utf8]{inputenc}
\usepackage{lmodern}
\usepackage{amsthm,amsfonts,amssymb,amsmath}

\newtheorem{claim}{Claim}
\newtheorem{definition}{Definition}

\title{CDF measures}
\author{Aditya Padiyar}

\begin{document}

\maketitle

\begin{claim}
  The collection of finite unions of (possibly unbounded) intervals of the type $(a,b]$ forms an algebra $\mathcal{A}$ of subsets of $\mathbb{R}$
\end{claim}

\begin{claim}
  Every $A \in \mathcal{A}$ has finitely many connected components
\end{claim}

$F,G:\mathbb{R} \to \mathbb{R}$ denote nondecreasing, right continuous functions

\begin{definition}
  For an interval $I$ with endpoints $a$ and $b$ we define \[\text{length}_F(I)=F(b)-F(a)\] where we use the appropriate limits of $F$ whenever the interval is unbounded
\end{definition}

\begin{claim}
  Define $\mu_F:\mathcal{A} \to [0,\infty]$ by \[\mu_F(A)=\sum_{\text{connected component } I } \text{length}_F(I)\] 
  Then $\mu_F$ is a measure on the algebra $\mathcal{A}$
\end{claim}

We now use the Caratheodory extension theorem to extend to $\mathcal{B}(\mathbb{R})$

\begin{claim}
  Suppose $I=(a,b]$ is bounded. Then we have \[\int_I F d\mu_G + \int_I G d \mu_F=F(b)G(b)-F(a)G(a)\] assuming $F$ and $G$ have no common points of discontinuity
\end{claim}

\begin{proof}
  Let $D$ be the diagonal of the square $I \times I$
  
  \begin{itemize}
    \item Let $S_G$ denote the set of points of discontinuity of $G$. A standard result tells us that $S_G$ is countable. Since $F$ and $G$ have no common discontinuities we have $\mu_F(S_G)=0$
    \item Now $\mu_D=\int_I \mu_G\{x\} d\mu_F=\int_I 1_{S_G} d\mu_F = 0$
    \item We have $F(x)=\mu_F(a,x] + F(a)$ and $G(x)=\mu_G(a,x]+G(a)$
  \end{itemize}
  
  Use Tonelli's theorem to rewrite the left hand side as 
  \[F(a)\text{length}_G(I)+G(a)\text{length}_F(I)+\text{length}_F(I)\text{length}_G(I)+\mu(D)\]
  Now simplify the equation to get the desired result
\end{proof}

\end{document}
